% ============================================================
%  System Architecture section --- ModSim workshop paper.
%  Placement: immediately after Related Work.
%  Target: IEEE conference template (IEEEtran, two-column).
%
%  >>> PREAMBLE REQUIREMENTS (add once, if not already present):
%     \usepackage{tikz}
%     \usetikzlibrary{positioning,arrows.meta,fit,backgrounds}
%     \usepackage{amsmath}          % \mathrm, aligned math in text
%     % (booktabs/amssymb/xcolor are already loaded for the related-work table)
%
%  The figure is drawn in native TikZ at \columnwidth, so its fonts match the
%  body text and it needs no external asset. It floats to the top of a column
%  ([t]); \label{fig:architecture} is referenced from the opening paragraph.
% ============================================================

\section{System Architecture}
\label{sec:architecture}

ModSim rests on a single architectural commitment: \emph{separate modular
meaning from physics, and make the boundary between them a narrow, explicit
contract.} The core decides \emph{whether} two connectors may mate and \emph{what
it means} when they do; a backend decides only whether the requested constraint
physically exists. A backend is never asked ``is this a valid dock?''---it is
asked ``create this constraint,'' and it answers yes or no.
Figure~\ref{fig:architecture} shows the resulting stack, in which URDF is an
imported mechanical asset, Robot Pack YAML the semantic description,
\texttt{WorldState} the canonical runtime state, and topology graphs
\emph{derived} views. The rest of this section develops that split and the
verification it enables.

% ---------------------------------------------------------------
%  ARCHITECTURE FIGURE  (single column, native TikZ)
% ---------------------------------------------------------------
\begin{figure}[t]
  \centering
  \resizebox{\columnwidth}{!}{% guarantee the figure fits one IEEE column
  \begin{tikzpicture}[
      font=\footnotesize,
      >={Stealth[length=5pt]},
      box/.style={draw, rounded corners=2pt, align=center,
                  inner sep=3pt, minimum height=6mm},
      input/.style={box, fill=gray!8, minimum width=16mm},
      core/.style={box, fill=ourrow, minimum width=14mm, minimum height=8mm},
      backend/.style={box, fill=gray!5, minimum width=30mm, minimum height=8mm},
      lbl/.style={font=\scriptsize\itshape, text=gray!55},
    ]
    % --- core row (defined first; inputs and the title position off this box) ---
    \node[core] (world) {\texttt{WorldState}\\{\scriptsize(canonical)}};
    \node[core, right=2mm of world] (dock) {Docking\\engine};
    \node[core, right=2mm of dock] (log) {Event log\\{\scriptsize+ metrics}};

    % dashed container = the backend-agnostic core (title added separately, below)
    \begin{scope}[on background layer]
      \node[draw, dashed, rounded corners=3pt, fit=(world)(dock)(log),
            inner sep=3.5mm] (corebox) {};
    \end{scope}

    % --- inputs, raised above the box and set over the outer core cells ---
    \node[input] (pack)  at ([yshift=13mm]world.north)
      {Robot Pack\\{\scriptsize(URDF + semantics)}};
    \node[input] (scene) at ([yshift=13mm]log.north) {Scene};

    % --- adapter contract (the boundary) ---
    \node[box, fill=black!78, text=white, below=6mm of corebox.south,
          minimum width=64mm, text width=60mm]
      (contract)
      {\textbf{Adapter contract}\\[1pt]{\scriptsize\ttfamily capabilities
       $\cdot$ load $\cdot$ step $\cdot$ snapshot\\
       create/remove\_connection $\cdot$ shutdown}};

    % --- backends (bottom row) ---
    \node[backend, below=4.5mm of contract.south west, anchor=north west]
      (mock) {\textbf{mock}\\{\scriptsize kinematic reference}};
    \node[backend, below=4.5mm of contract.south east, anchor=north east]
      (mj) {\textbf{MuJoCo}\\{\scriptsize weld-pool bridge}};

    % --- arrows: inputs feed the core, entering either side of the title ---
    \draw[->] (pack.south)  -- (corebox.north -| pack.south);
    \draw[->] (scene.south) -- (corebox.north -| scene.south);

    % title as a fieldset-style legend centred on the top border (drawn last; masks dashes)
    \node[fill=white, inner sep=1.5pt, font=\scriptsize\bfseries]
      at (corebox.north) {Backend-agnostic core};

    % core <-> contract : two labelled directions
    \draw[->] ([xshift=-12mm]corebox.south) -- ([xshift=-12mm]contract.north)
      node[lbl, midway, left=1pt, align=right] {create /\\remove};
    \draw[<-] ([xshift=12mm]corebox.south) -- ([xshift=12mm]contract.north)
      node[lbl, midway, right=1pt, align=left] {snapshot};

    % contract <-> backends
    \draw[<->] ([xshift=15mm]contract.south west) -- (mock.north);
    \draw[<->] ([xshift=-15mm]contract.south east) -- (mj.north);

    % conformance oracle link between the two backends
    \draw[<->, dashed, gray!60] (mock.east) -- (mj.west)
      node[lbl, midway, above=0pt, fill=white, inner sep=1pt] {conformance};
  \end{tikzpicture}%
  }
  \caption{ModSim architecture. A portable Robot Pack (URDF geometry plus
    modular semantics) and a scene feed a backend-agnostic core that owns
    canonical \texttt{WorldState}, the docking engine, and the event log. A
    seven-method adapter contract is the sole boundary; the core commits a
    logical connection only after a backend confirms the constraint. The
    dependency-free kinematic \emph{mock} serves as reference semantics:
    cross-backend conformance asserts that mock and MuJoCo agree on connector
    frames, docking verdicts, and the committed event sequence.}
  \label{fig:architecture}
\end{figure}

\noindent\textbf{Portable modules over a narrow contract.}\ A \emph{Robot Pack} describes one module as split YAML layered over a standard
URDF: URDF supplies geometry and kinematics, and the pack adds what a physics
engine has no vocabulary for---named connectors and their frames, gender and
bidirectional compatibility, acceptance tolerances, and allowed mating
orientations---so one description drives every backend and a pack authored once
runs unchanged on the kinematic mock and on MuJoCo. The core reaches a backend
through only seven methods (\texttt{capabilities}, \texttt{load}, \texttt{step},
\texttt{snapshot}, \texttt{create\_physical\_connection},
\texttt{remove\_physical\_connection}, \texttt{shutdown}), with engine-specific
settings such as gravity and timestep on the constructor rather than in the
shared interface. \texttt{capabilities()} is load-bearing, not documentation: the
docking engine checks \texttt{supports\_runtime\_constraints} before a commit, so
a backend that cannot enforce a connection yields an explicit \texttt{DockFailed}
rather than a silently faked latch. Registration is open and lazy---importing the
core never imports MuJoCo~\cite{todorov2012mujoco}---so adding an engine such as
Isaac Sim~\cite{nvidia_isaacsim} means implementing seven methods, not touching
the core.

\noindent\textbf{An event-sourced core and a weld-pool bridge.}\ The core owns a canonical \texttt{WorldState} mutated only by applying events, so
the log is a complete, replayable account and assembly identity is derived from
membership
(\(\mathrm{AssemblyId}=\text{``assembly:''}+\min(\text{member ids})\)), yielding
the same identifier across runs and replays---the determinism that log-replaying
reconfiguration experiments require~\cite{liu2019distributed}. A logical
connection is committed only through a two-phase commit---re-check acceptance
against the current snapshot, request the constraint, emit \texttt{DockCommitted}
only after the backend confirms it, or \texttt{DockFailed} on refusal---and
within a step releases precede detection, so a connection cannot break and
re-form at once; there is thus one source of truth for connectivity, never ahead
of the physics. Realising that contract on MuJoCo, whose topology is fixed at
compile time~\cite{todorov2012mujoco}, means a dock cannot \emph{create} a
constraint: the adapter reserves a pool of inactive welds, and a dock claims,
re-points, and activates one while an exhausted pool refuses rather than fakes a
latch. Because semantics live between connector frames while a weld constrains
bodies, the conversion
\(T_{\mathrm{bodyA}\leftarrow\mathrm{bodyB}}
= T_{\mathrm{bodyA}\leftarrow\mathrm{connA}}\,
T_{\mathrm{connA}\leftarrow\mathrm{connB}}\,
T_{\mathrm{connB}\leftarrow\mathrm{bodyB}}\)
is isolated in one module, unit-testable without a running simulation, with
articulated-child connectors mapped back to a root free-joint pose before the
weld activates; where the engine materialises connector frames natively (MuJoCo
sites), the core prefers those measured frames so core and engine cannot disagree
about where a connector is.

\noindent\textbf{Verification by cross-backend agreement.}\ Separating semantics from physics does more than keep the code clean: it makes
correctness checkable by differential testing. The mock backend is deliberately kinematic---it models
only that a connection makes two modules move as one rigid body and that removing
it lets them move independently---and stays permanently as both the
continuous-integration backend and the reference semantics. A cross-backend
conformance suite runs the same pack and scene through every installed backend
and asserts agreement on the semantic picture: the same modules and assemblies,
connector world frames, and \emph{committed event sequence}, with identical
connection identities and mating orientation indices. Comparisons run with
gravity disabled, isolating adapter correctness from dynamics: poses legitimately
diverge once physics acts, but the \emph{decisions} must not, since each is made
by backend-agnostic core code from the reported state. A disagreement therefore
points at a mismapped link or a mis-framed connector, not at modules that merely
fell---turning ``did this simulate correctly?'' into a concrete, automatable
assertion between an independent kinematic model and a full-physics one.
